\begin{definition}[Generalized central factorial numbers of the second kind]
  For integers $0 \leq k \leq m$, and an arbitrary integer $t$, the generalized
  central factorial numbers of the second kind are defined by the relation
  \begin{align*}
    n^m = \sum_{k=0}^{\infty} T_t(m,k) \cdot \centralFactorial{(n-t)}{k}.
  \end{align*}
\end{definition}
It is also straightforward that
\begin{align*}
  T_t (m,k) = \frac{\delta^{k} t^m}{k!},
\end{align*}
by Newton's
formula~\eqref{lem:central-newtons-formula-for-powers}.
For example,
\begin{table}[H]
  \begin{center}
    \setlength\extrarowheight{-6pt}
    \begin{tabular}{c|ccccccccc}
      $n/k$ & 0 & 1 & 2 & 3 & 4 & 5 & 6 & 7 & 8 \\ [3px]
      \hline
      0 & 1 &   &   &   &   &   &   &   &   \\
      1 & 0 & 1 &   &   &   &   &   &   &   \\
      2 & 0 & 0 & 1 &   &   &   &   &   &   \\
      3 & 0 & $\tfrac{1}{4}$ & 0 & 1 &   &   &   &   &   \\
      4 & 0 & 0 & 1 & 0 & 1 &   &   &   &   \\
      5 & 0 & $\tfrac{1}{16}$ & 0 & $\tfrac{5}{2}$ & 0 & 1 &   &   &   \\
      6 & 0 & 0 & 1 & 0 & 5 & 0 & 1 &   &   \\
      7 & 0 & $\tfrac{1}{64}$ & 0 & $\tfrac{91}{16}$ & 0 & $\tfrac{35}{4}$ & 0 & 1 &   \\
      8 & 0 & 0 & 1 & 0 & 21 & 0 & 14 & 0 & 1 \\
    \end{tabular}
  \end{center}
  \caption{Values of generalized central factorial numbers $T_0(n,k)$.
    See also the sequences \href{https://oeis.org/A000000}{\texttt{A000000}},
    and \href{https://oeis.org/A000000}{\texttt{A000000}} in~\cite{sloane2003line}.
  }
\end{table}


\begin{table}[H]
  \begin{center}
    \setlength\extrarowheight{-6pt}
    \begin{tabular}{c|ccccccccc}
      $n/k$ & 0 & 1 & 2 & 3 & 4 & 5 & 6 & 7 & 8 \\ [3px]
      \hline
      0 & 1 &   &   &   &   &   &   &   &   \\
      1 & 1 & 1 &   &   &   &   &   &   &   \\
      2 & 1 & 2 & 1 &   &   &   &   &   &   \\
      3 & 1 & $\tfrac{13}{4}$ & 3 & 1 &   &   &   &   &   \\
      4 & 1 & 5 & 7 & 4 & 1 &   &   &   &   \\
      5 & 1 & $\tfrac{121}{16}$ & 15 & $\tfrac{25}{2}$ & 5 & 1 &   &   &   \\
      6 & 1 & $\tfrac{91}{8}$ & 31 & 35 & 20 & 6 & 1 &   &   \\
      7 & 1 & $\tfrac{1093}{64}$ & 63 & $\tfrac{1491}{16}$ & 70 & $\tfrac{119}{4}$ & 7 & 1 &   \\
      8 & 1 & $\tfrac{205}{8}$ & 127 & $\tfrac{483}{2}$ & 231 & 126 & 42 & 8 & 1 \\
    \end{tabular}
  \end{center}
  \caption{
    Values of generalized central factorial numbers $T_1(n,k)$.
    See also the sequences \href{https://oeis.org/A395860}{\texttt{A395860}},
    and \href{https://oeis.org/A395861}{\texttt{A395861}} in the OEIS~\cite{sloane2003line}.
  }
  \label{tab:central-factorial-numbers-t1}
\end{table}


\begin{table}[H]
  \begin{center}
    \setlength\extrarowheight{-6pt}
    \begin{tabular}{c|ccccccccc}
      $n/k$ & 0 & 1 & 2 & 3 & 4 & 5 & 6 & 7 & 8 \\ [3px]
      \hline
      0 & 1 &   &   &   &   &   &   &   &   \\
      1 & 2 & 1 &   &   &   &   &   &   &   \\
      2 & 4 & 4 & 1 &   &   &   &   &   &   \\
      3 & 8 & $\tfrac{49}{4}$ & 6 & 1 &   &   &   &   &   \\
      4 & 16 & 34 & 25 & 8 & 1 &   &   &   &   \\
      5 & 32 & $\tfrac{1441}{16}$ & 90 & $\tfrac{85}{2}$ & 10 & 1 &   &   &   \\
      6 & 64 & $\tfrac{931}{4}$ & 301 & 190 & 65 & 12 & 1 &   &   \\
      7 & 128 & $\tfrac{37969}{64}$ & 966 & $\tfrac{12411}{16}$ & 350 & $\tfrac{371}{4}$ & 14 & 1 &   \\
      8 & 256 & $\tfrac{6001}{4}$ & 3025 & 3003 & 1701 & 588 & 126 & 16 & 1 \\
    \end{tabular}
  \end{center}
  \caption{Values of generalized central factorial numbers $T_2(n,k)$.
    See also the sequences \href{https://oeis.org/A000000}{\texttt{A000000}},
    and \href{https://oeis.org/A000000}{\texttt{A000000}} in~\cite{sloane2003line}.
  }
\end{table}


\begin{table}[H]
  \begin{center}
    \setlength\extrarowheight{-6pt}
    \begin{tabular}{c|ccccccccc}
      $n/k$ & 0 & 1 & 2 & 3 & 4 & 5 & 6 & 7 & 8 \\ [3px]
      \hline
      0 & 1 &   &   &   &   &   &   &   &   \\
      1 & 3 & 1 &   &   &   &   &   &   &   \\
      2 & 9 & 6 & 1 &   &   &   &   &   &   \\
      3 & 27 & $\tfrac{109}{4}$ & 9 & 1 &   &   &   &   &   \\
      4 & 81 & 111 & 55 & 12 & 1 &   &   &   &   \\
      5 & 243 & $\tfrac{6841}{16}$ & 285 & $\tfrac{185}{2}$ & 15 & 1 &   &   &   \\
      6 & 729 & $\tfrac{12753}{8}$ & 1351 & 585 & 140 & 18 & 1 &   &   \\
      7 & 2187 & $\tfrac{372709}{64}$ & 6069 & $\tfrac{53011}{16}$ & 1050 & $\tfrac{791}{4}$ & 21 & 1 &   \\
      8 & 6561 & $\tfrac{167943}{8}$ & 26335 & $\tfrac{35049}{2}$ & 6951 & 1722 & 266 & 24 & 1 \\
    \end{tabular}
  \end{center}
  \caption{Values of generalized central factorial numbers $T_3(n,k)$.
    See also the sequences \href{https://oeis.org/A000000}{\texttt{A000000}},
    and \href{https://oeis.org/A000000}{\texttt{A000000}} in the OEIS~\cite{sloane2003line}.
  }
  \label{tab:central-factorial-numbers-t3}
\end{table}


\input{sections/tables/01_fig_oeis_central_factorial_numbers_2n_2k.tex}

Thus, we have the formula for sums of powers in terms of
generalized central factorial numbers of the second kind

\begin{proposition}[Generalized Central factorial sums]
  \label{prop:generalized-central-factorial-sums}
  For integers $n,m \geq 0$, and an arbitrary integer $t$,
  \begin{align*}
    \KnuthRFoldSum{1}{n}{m} = \frac{1}{2} \sum_{k=0}^{m}
    \left[
      \binom{n -t + \frac{k}{2} +1}{k+1} - \binom{1 -t + \tfrac{k}{2}}{k+1}
      + \binom{n -t + \frac{k}{2}}{k+1} - \binom{-t + \tfrac{k}{2}}{k+1}
    \right] k! T_t (m,k).
  \end{align*}
\end{proposition}
